\renewcommand{\Anno}{2000}

\renewcommand{\NumeroProblema}{3}

\renewcommand{\Testo}{Per poter usufruire di una borsa di studio, uno
studente deve superare un test costituito da 10 domande.  Per ogni domanda il punteggio è $+1$ se lo studente risponde correttamente, $-1$ se lo studente risponde in modo errato, $0$ se lo studente non risponde.  Il test si considera superato se il punteggio totale, ottenuto sommando i punteggi ottenuti nelle 10 domande, è maggiore od uguale a $5$.
	
Diciamo che uno studente ha preparazione $p$ (con $0\leq p\leq 1$) se, per ogni domanda a cui prova a rispondere, egli ha probabilità $p$ di fornire la risposta giusta.
	
Determinare, in funzione di $p$, la strategia che lo studente deve seguire (cioè il numero di domande alle quali provare a rispondere) per massimizzare la sua probabilità di superare il test.  Calcolare anche tale probabilità (sempre in funzione di
$p$).}

\renewcommand{\Risposta}{La probabilità di superare il test, in funzione
	di $p$, è data da
	$$\mbox{Probabilit\`{a} }(p)=\left\{
	\renewcommand{\arraystretch}{1.4}
	\begin{array}{l@{\hspace{10mm}}l}
		p^{5} & \mbox{se }0\leq p\leq 1/6,  \\
		7p^{6}-6p^{7} & \mbox{se }1/6\leq p\leq 1/4,  \\
		28p^{9}-63p^{8}+36p^{7} & \mbox{se }1/4\leq p\leq 1.
	\end{array}
	\right.$$
	
Per quanto riguarda la strategia:
\begin{itemize}
	\item  se $p=0$ tutte le strategie sono equivalenti;
	
	\item  se $0<p<1/6$ conviene rispondere a 5 domande;
	
	\item se $p=1/6$ conviene indifferentemente rispondere a 5 o 7 domande;
	
	\item  se $1/6<p<1/4$ conviene rispondere a 7 domande;
	
	\item se $p=1/4$ conviene indifferentemente rispondere a 7 o 9 domande;
	
	\item se $1/4<p<1$ conviene rispondere a 9 domande; 
			
	\item se $p=1$ è indifferente rispondere ad un qualunque numero maggiore od uguale a 5 di domande.
\end{itemize} 
}

\renewcommand{\Aiutino}{Dato un intero $n$, maggiore od uguale a 5 e
minore od uguale a 10, indichiamo con $f_{n}(p)$ la probabilità di superare il test rispondendo ad esattamente $n$ domande ed avendo preparazione $p$.  Il problema chiede di determinare, al variare di $p$, il massimo tra tutte le $f_{n}$.  Con qualche calcolo si ottiene che tali funzioni sono
\begin{eqnarray*}
	f_{5}(p) & = & p^{5},  \\
	f_{6}(p) & = & p^{6},  \\
	f_{7}(p) & = & p^{7}+7p^{6}(1-p),  \\
	f_{8}(p) & = & p^{8}+8p^{7}(1-p),  \\
	f_{9}(p) & = & p^{9}+9p^{8}(1-p)+36p^{7}(1-p)^{2},  \\
	f_{10}(p) & = & p^{10}+10p^{9}(1-p)+45p^{8}(1-p)^{2}.
\end{eqnarray*}

Si verifica ora con semplici disequazioni che $f_{6}\leq f_{5}$, $f_{8}\leq f_{7}$ e $f_{10}\leq f_{9}$, con uguaglianza se e solo se $p=0$ oppure $p=1$. Ci si riduce cos\`{\i} a determinare il massimo
tra le $f_{n}$ di indice dispari.}

\renewcommand{\SoluzioneUfficiale}{
}

\renewcommand{\SoluzioniAlternative}{
}

\renewcommand{\Osservazioni}{
}

\renewcommand{\Commenti}{
}
% Aggiunta redazionale: la soluzione preesistente non viene modificata.
\providecommand{\SoluzioneRedazionale}{}
\renewcommand{\SoluzioneRedazionale}{
% SSSUP-BOZZA-AI: SoluzioneRedazionale; autore=ChatGPT; data=2026-11-08; stato=DA-VERIFICARE
% SSSUP-PROVENIENZA: AI: elaborazione nuova dal testo
Assumiamo che gli esiti delle risposte alle diverse domande siano indipendenti, ciascuno con probabilità $p$ di successo. La sola probabilità marginale indicata nel testo, senza questa ipotesi, non determina la probabilità di superare il test.

Rispondendo a $n$ domande e indovinandone $X$, il punteggio è $2X-n$. Per $n<5$ non si può superare il test. Inoltre rispondere a $2j$ domande non è meglio che rispondere alle prime $2j-1$, per $2j=6,8,10$: ogni esito che supera la soglia con $2j$ risposte la supera già dopo le prime $2j-1$. Per $0<p<1$ il confronto è stretto.

Restano le probabilità
\[
f_5=p^5,\qquad f_7=p^6(7-6p),\qquad
f_9=p^7(36-63p+28p^2).
\]
Esse si confrontano senza risolvere equazioni di grado elevato, perché
\[
f_7-f_5=p^5(1-p)(6p-1),\qquad
f_9-f_7=7p^6(1-p)^2(4p-1).
\]
Per $0<p<1/6$ vince $f_5$; per $1/6<p<1/4$ vince $f_7$; per $1/4<p<1$ vince $f_9$. Nei due valori di confine sono ottimali entrambe le strategie adiacenti. Per $p=0$ nessuna strategia supera il test; per $p=1$ basta rispondere ad almeno cinque domande.

Il massimo è dunque
\[
\begin{cases}
p^5,&0\leq p\leq1/6,\\
7p^6-6p^7,&1/6\leq p\leq1/4,\\
36p^7-63p^8+28p^9,&1/4\leq p\leq1.
\end{cases}
\]
}
